% Load the beamer-sharcnet document class...
\documentclass[tabu]{beamer-uwindsor}

% Have latex inform when another run is needed...
\usepackage{rerunfilecheck}

% Use the csquotes package...
\usepackage{csquotes}

% Use the listings package...
\usepackage{listings}

\lstloadlanguages{Haskell}

\lstnewenvironment{code}
{\lstset{}%
	\csname lst@SetFirstLabel\endcsname}
{\csname lst@SaveFirstLabel\endcsname}
\lstset{
	basicstyle={\scriptsize\ttfamily\singlespacing},
	flexiblecolumns=false,
	basewidth={0.5em,0.45em},
	literate={+}{{$+$}}1 {/}{{$/$}}1 {*}{{$*$}}1 {=}{{$=$}}1
	{>}{{$>$}}1 {<}{{$<$}}1 {\\}{{$\lambda$}}1
	{\\\\}{{\char`\\\char`\\}}1
	{->}{{$\rightarrow$}}2 {>=}{{$\geq$}}2 {<-}{{$\leftarrow$}}2
	{<=}{{$\leq$}}2 {=>}{{$\Rightarrow$}}2
	{\ .\ }{{$\circ$}}2
	{>>}{{>>}}2 {>>=}{{>>=}}2
	{|}{{$\mid$}}1
}

\lstnewenvironment{spec}
{\lstset{}%
	\csname lst@SetFirstLabel\endcsname}
{\csname lst@SaveFirstLabel\endcsname}
\lstset{
	basicstyle={\scriptsize\ttfamily\singlespacing},
	flexiblecolumns=false,
	basewidth={0.5em,0.45em},
	literate={+}{{$+$}}1 {/}{{$/$}}1 {*}{{$*$}}1 {=}{{$=$}}1
	{>}{{$>$}}1 {<}{{$<$}}1 {\\}{{$\lambda$}}1
	{\\\\}{{\char`\\\char`\\}}1
	{->}{{$\rightarrow$}}2 {>=}{{$\geq$}}2 {<-}{{$\leftarrow$}}2
	{<=}{{$\leq$}}2 {=>}{{$\Rightarrow$}}2
	{\ .\ }{{$\circ$}}2
	{>>}{{>>}}2 {>>=}{{>>=}}2
	{|}{{$\mid$}}1
}

%
% Use TrueType fonts (i.e., xelatex).
%
% Use "Charis SIL" font even though it is not sans-serif as
% it is easy-to-read, has all Unicode characters, is under
% the open SIL font licence, and is free.
%
% As SIL does not have a monospaced font, Liberation Mono
% is used instead.
%
\usepackage{xltxtra,fontspec,xunicode}
\defaultfontfeatures{Scale=MatchLowercase}
\setromanfont{Charis SIL}
\setsansfont{Charis SIL}
\setmainfont{Charis SIL}
%\setmonofont[Scale=0.70]{Liberation Mono}
\setmonofont{Liberation Mono}

%
% Define slide show and presenter-specific globals...
%
\providecommand{\basetexfilename}{slides}
\providecommand{\emailshane}{peelar@uwindsor.ca}
\title[Masters Thesis Defence]{Accommodating prepositional phrases in a highly modular
	natural language query interface to semantic web triplestores using a novel event-based denotational semantics for English and a set of parser combinators.}
\date{December 12th 2016}
\author[Shane Peelar]{%
  Shane Peelar, BSc (Hons) \texorpdfstring{\\}{}%
  {\footnotesize\href{mailto:\emailshane}{\emailshane}} \texorpdfstring{\\}{}%
}
\institute[University of Windsor]{%
  University of Windsor, \texorpdfstring{\\}{}%
  Windsor, Ontario, Canada \texorpdfstring{\\}{}%
  Copyright \copyright{} 2016 \insertshortauthor{}. All Rights Reserved. \texorpdfstring{\\}{}%
  ~ \texorpdfstring{\\}{}%
}
\subject{Intro}

\usepackage{hyperref}

\usepackage{xcolor}
\hypersetup{
	colorlinks,
	linkcolor={red!50!black},
	citecolor={blue!50!black},
	urlcolor={blue!80!black}
}

% bibliography...
\usepackage[backend=biber,bibstyle=numeric-comp,sorting=ydnt,natbib=true,mcite=true,maxnames=100,url=true,isbn=false,doi=false,uniquename=init,giveninits=true,hyperref=true,backref=true,date=edtf,sortcites]{biblatex}
\renewcommand{\subtitlepunct}{\addcolon\addspace}

%\hypersetup{pdfpagemode=FullScreen}


% Load support for hypertext references...
\usepackage{url}

% Turn off the navigation bar...
\beamertemplatenavigationsymbolsempty

% Support arithmetic with numbers...
\usepackage{fp}

% Use Tikz
\usepackage{tikz}

%
% Make it easy to output a number in roman numeral form...
%
\makeatletter
\newcommand*{\asroman}[1]{\expandafter\@slowromancap\romannumeral #1@}
\makeatother

%
% Redefine the part page to subtract one from the part # since the first
% part of this document represents the document title page and overview
% preamble...
%
\setbeamertemplate{part page}
{
  \begin{centering}
    \FPeval{\result}{clip(\insertpartnumber-1)}
    {\usebeamerfont{part name}\usebeamercolor[fg]{part name}\partname~\asroman{\result}}
    \vskip1em\par
    \begin{beamercolorbox}[sep=16pt,center,shadow=false,rounded=true]{part title}
      \usebeamerfont{part title}\insertpart\par
    \end{beamercolorbox}
  \end{centering}
}

%
% Define how beamer shows auto-generated pages...
%
\makeatletter
\newcommand*{\resetsectioncount}{
  \beamer@tocsectionnumber=0\relax%
  \setcounter{section}{0}%
}
\makeatother

\makeatletter
\AtBeginPart{%
  % Restart section numbering within each part...
  \resetsectioncount%
  % And output (if applicable) the part page...
  \ifthenelse{\boolean{showpartpage}}{%
    \frame{\partpage}%
  }{}%
}
\makeatother

\AtBeginSection[]{%
  \ifthenelse{\boolean{showsectiontoc}}{%
    \begin{frame}{Table of Contents}%
      \tableofcontents[sectionstyle=show/hide,subsectionstyle=show/show/hide]%
    \end{frame}%
  }{}%
}

% Bibliography...
\addbibresource{thesis.bib}

% Title Page Logo
\titlegraphic{\includegraphics[width=2cm]{logos/uwin-logo-horz-3col.eps}}

% 
% The presentation slide content follows...
%
\begin{document}
	% Don't show logo or footer on title page...
	{
	\setbeamertemplate{logo}{}
	\setbeamertemplate{footline}{} 
	\begin{frame}[label=TitleSlide,noframenumbering]
	  \titlepage
	\end{frame}
	}
	
	\begin{frame}{Thesis Committee}
		\begin{itemize}
			\item Supervisors: Dr. Richard Frost and Dr. Robert Kent
			\item Internal Reader: Dr. Luis Rueda
			\item External Reader: Dr. Richard Caron
			\item Chair: Dr. Alioune Ngom
		\end{itemize}
	\end{frame}

	\begin{frame}{Outline}
		\small
		\tableofcontents
	\end{frame}
	
	\sectionwithouttoc{Introduction}

	
	\subsectionwithouttoc{Motivation}
	\begin{frame}{Motivation}
		\begin{itemize}
			\item Internet of Things
			\begin{itemize}
				\item Self driving cars, coffee machines, refrigerators, ovens...
				\item Able to be operated remotely through {\em alternative interfaces}
				\item Accessibility
				\item Natural Language Interfaces
			\end{itemize}
			\item The Semantic Web
		\end{itemize}
	\end{frame}

	\subsectionwithouttoc{The Semantic Web}
	\begin{frame}{The Semantic Web}
		\begin{itemize}
			\item The Semantic Web (part of Web 3.0)
			\begin{itemize}
				\item Distributed and decentralized network of databases
				\item Store facts in the form of {\em triples}
				\item Resource Description Framework (RDF) \cite{w3csemanticweb}
				\item SPARQL Protocol and RDF Query Language (SPARQL)
			\end{itemize}
			\item Can be thought of, informally, as applying the model of the World Wide Web to databases
			\item SPARQL is low-level (like SQL)
		\end{itemize}
	\end{frame}
	
	%\subsectionwithouttoc{Resource Description Framework (RDF)}
	\begin{frame}{Resource Description Framework (RDF)}
		\begin{itemize}
			\item Data is organized into {\em triples}
			\item A {\em triple} is a 3-tuple that has the form \texttt{(subject, predicate, object)}
			\item \texttt{subject}, \texttt{predicate}, and \texttt{object} are Uniform Resource Identifiers (URIs)\cite{w3csemanticweb}.  A URI may denote a {\em name}, {\em location}, or possibly both at once.  Examples:
			\begin{itemize}
				\item URL
				\item ISBN of a book
				\item Telephone number
			\end{itemize}
			\item A database containing these triples is commonly called a {\em triplestore}
			\item Example triple: \texttt{(jane, purchased, pencil\_1)}
			\item Triplestores that describe facts in terms of {\em entities} are called {\em entity-based triplestores}
		\end{itemize}
	\end{frame}
	
	%\subsectionwithouttoc{SPARQL}
	\defverbatim[colored]\lstI{
		\begin{lstlisting}[basicstyle=\ttfamily,keywordstyle=\color{red}]
	PREFIX foaf:  <http://xmlns.com/foaf/0.1/>
	SELECT ?name ?email
	FROM <http://www.w3.org/People/Berners-Lee/card>
	WHERE {
	{      SELECT DISTINCT ?person ?name WHERE { 
	?person foaf:name ?name 
	} ORDER BY ?name LIMIT 10 OFFSET 10    }
	}
		\end{lstlisting}
	}
	
	\begin{frame}{SPARQL Protocol and RDF Query Language}
		\begin{itemize}
			\item RDF Query Language
			\item Similar to SQL
			\item Queries are submitted to {\em SPARQL endpoints}
			\item Form patterns that define restrictions on the elements of triples
			\item Similar to pattern matching in functional programming
			\item Low-level
			\item Example: \lstI
		\end{itemize}
	\end{frame}
	
	\subsectionwithouttoc{The Problem}
	\begin{frame}[allowframebreaks]{The Problem}
		\begin{itemize}
			\item SPARQL is low level and not user-friendly
			\item Natural Language Interfaces
				\begin{itemize}
					\item Require little technical knowledge to use
					\item Require minimal effort on part of the user
					\item Improve accessibility
					\item Area of active research, with several previous attempts (non-trivial)%TODO: HERE
				\end{itemize}
		\end{itemize}
	\end{frame}

	\subsectionwithouttoc{Previous Work}
	\begin{frame}[allowframebreaks]{Previous Work}
		\begin{itemize}
			\item ORAKEL (2007) \cite{cimiano2007orakel}
			\begin{itemize}
				\item Ontology-aware
				\item English interface
				\item Based on Montague semantics
				\item Supports quantification, negation, conjunction
				\item Attempts to directly convert input query into a SPARQL query
			\end{itemize}
			\item QuestIO (2008) \cite{tablan2008natural}
			\begin{itemize}
				\item Ontology-aware
				\item English interface
				\item Keyword oriented
				\item Attempts to match keywords against concepts to hone down queries
				\item Uses SeRQL, similar to SPARQL, to form queries
				\item Sentences transformed into SeRQL queries using formal semantics
			\end{itemize}
		
			\pagebreak
			\item AutoSPARQL (2011) \cite{lehmann2011autosparql}
			\begin{itemize}
				\item Supervised machine learning
				\item Constructs query trees using keywords
				\item Query trees are converted into SPARQL queries
				\item Feedback oriented
				\begin{itemize}
				\item User selects candidates from returned query set that best matches what the user is looking for
				\item User is expected to be actively involved in refining subsequent results
				\end{itemize}
			\end{itemize}
		\end{itemize}
	\end{frame}
	
	\subsectionwithouttoc{Shortcomings of Previous Work}
	\begin{frame}{Shortcomings of Previous Work}
		
		%\textbf{The Problem}:
		Previous work has seen success, but the English NLIs have a shortcoming in that they do not allow for prepositional phrases in queries, and hence have a limited coverage of the English language.
		%\newline
		
	\end{frame}

	\subsectionwithouttoc{New Approach}
	\begin{frame}{New Approach}
		
		The work presented in this thesis draws on two main concepts:
		\begin{itemize}
			\item Executable Attribute Grammars\cite{frosthafiz2008}
			\item Event-Based Denotational Semantics\cite{frostagboola2014}.
		\end{itemize}
	\end{frame}
	
	\subsubsectionwithouttoc{Executable Attribute Grammars}
	\begin{frame}{Executable Attribute Grammars}
		\begin{itemize}
		\item A natural way to process Natural Language queries \cite{frosthafiz2008}
		\item Allow top-down rather than bottom-up parsing, so are highly modular \cite{frosthafiz2008}
		\item Efficient top-down parsing has been demonstrated previously \cite{frosthafiz2008}
		\item In the parser we used:
		\begin{itemize}
			\item Left-recursive grammars are supported directly \cite{frosthafiz2008}
			\item Ambiguity in grammars is tracked, producing parse trees for all possible interpretations of a grammar \cite{frosthafiz2008}
		\end{itemize}
		%\item This allows ambiguous grammars to be used, simplifying development
		\item We modified the parser described in \cite{frosthafiz2008} to support using non-referentially transparent functions as attributes using the IO Monad
		\end{itemize}
	\end{frame}

	\begin{frame}{Functional Programming}
		\begin{itemize}
			\item Original parser written in Haskell \cite{frosthafiz2008}
			\item Referential transparency or ``pure'' functions
			\item Used ``memoization'' and lazy-evaluation to efficiently parse left recursive ambiguous grammars
			%\item Lazy-evaluation
			\item Extended to IO Monad to handle non-referentially transparent functions as attributes
		\end{itemize}
		
	\end{frame}
	
	\begin{frame}[allowframebreaks]{Denotational Semantics}
		Introduction to denotational semantics:
		%TODO: Explain quantification see frost2013event prepositions
		\begin{itemize}
			\item Montague Semantics
			\begin{itemize}
				\item Described by Dr. Richard Montague in 1970\cite{dowty2012introduction}
				\item Uses {\em characteristic functions} in {\em higher order logic} as denotations for English words %TODO:
				\item Not computationally tractable due to quantification: must traverse entire universe of discourse
				\item Denotations: %TODO
				\begin{itemize}
					\item Common nouns: $cnoun : \left(\forall ent\right) ent \rightarrow bool$
					\item Intransitive verbs: $intrans : \left(\forall ent\right) ent \rightarrow bool$
					\item Proper nouns: $pnoun : \left(\left(\forall ent\right) ent \rightarrow bool\right) \rightarrow bool$
				\end{itemize}
				\item ``phobos spins'':
				\begin{itemize}
					\item $phobos : \left(\left(\forall ent\right) ent \rightarrow bool\right) \rightarrow bool$
					\item $spins : \left(\forall ent\right) ent \rightarrow bool$
					\item $phobos\ spins \Rightarrow bool$
				\end{itemize}
			\end{itemize}
			\pagebreak
			\item FLMS
			\begin{itemize}
				\item Described by Frost et al. in 1989\cite{frost1989constructing}, based on Montague's approach
				\item Uses set-theory instead of higher order logic to describe denotations for English words
				\item Includes an explicit denotation for transitive verbs, missing from Montague's approach\cite{frost2013event}
				\item Computationally tractable\cite{frost1989constructing}
				\item Denotations: %TODO
				\begin{itemize}
					\item Common nouns: set of entities
					\item Intransitive verbs: set of entities
					\item Proper nouns: a function from sets of entities to booleans
					\item Transitive verbs: denoted using a binary relation
				\end{itemize}
				\item ``phobos spins'':
				\begin{itemize}
					\item phobos : $\{entities\} \rightarrow bool$
					\item spins : $\{entities\}$
					\item $phobos\ spins \Rightarrow bool$
				\end{itemize}
			\end{itemize}
		\end{itemize}
	\end{frame}

	\begin{frame}{FLMS and Transitive Verbs}
		\begin{itemize}
			\item Transitive verbs are denoted using binary relations
			\item Binary relations are converted into functions by ``collecting'' them
			\item $transvb : \left( \{entities\} \rightarrow bool \right) \rightarrow \{entities\}$
			\item $discover\ phobos \Rightarrow \{\left(hall, \{deimos, phobos\}\right)\}$
			\item Beginnings of notion of Function Defined By Relation (FDBR)
			\item Passive transitive verbs are created by inverting the binary relation
			
		\end{itemize}
	\end{frame}
	
	%\subsubsectionwithouttoc{Event-Based Denotational Semantics}
	\begin{frame}[allowframebreaks]{FLMS on the Semantic Web}
		\begin{itemize}
			\item FLMS can be directly implemented on entity-based triplestores
			\item Binary relation $rel = \{(a_1, b_1), (a_2, b_2), \dots\}$ represented with triples $(a_1, rel, b_1), (a_2, rel, b_2)\dots$
			\item No prepositions: how to get contextual information from a triple?
			\begin{itemize}
				\item Use quads instead of triples
				\item Reification
				\item {\em Event-Based Triplestores}
			\end{itemize}
		\end{itemize}	
	\end{frame}

	%TODO: Font
	\defverbatim[colored]\lstIoo{
	\begin{lstlisting}[basicstyle=\ttfamily,keywordstyle=\color{red}]
		(jane, purchased, pencil_1)
	\end{lstlisting}
	}

	\defverbatim[colored]\lstIo{
		\begin{lstlisting}[basicstyle=\ttfamily,keywordstyle=\color{red}]
		(event1, subject, jane)
		(event1, type, purchase)
		(event1, object, pencil)
		(event1, cost, 1)
		(event1, currency, CAD)
		\end{lstlisting}
	}

	%\sectionwithouttoc{Event-Based Triplestores}
	\begin{frame}{Event-Based Triplestores}
		\begin{itemize}
			\item Triples describe facts about events
			\item Information about entities and their relationships to one another may be gleaned from the events in which they occur
			\item Additional information about an event may be added by simply adding new triples
			\item Key motivation: directly supports reification on triples\cite{frostagboola2014}
		\end{itemize}
		Example (traditional):
		{\lstIoo}
		Example (event-based):
		{\lstIo}
	\end{frame}

	\subsubsectionwithouttoc{Event-Based Denotational Semantics}
	\begin{frame}[allowframebreaks]{Event-Based Denotational Semantics}
		\begin{itemize}
			\item EV-FLMS 
			\begin{itemize}
				\item Described by Frost et al. in 2013\cite{frost2013event}, based on FLMS
				\item Intended for use with {\em event-based triplestores}
				\item Includes a denotation for prepositional phrases
				\item Key difference with existing approaches: direct translation to SPARQL is not performed
				\item Triplestore access abstracted through \texttt{getts} functions
				\begin{itemize}
					\item $get\_members\ set = getts(ANY, type, membership) \cap getts(ANY, object, set)$
					\item $spins = get\_members\ spin$
				\end{itemize}
				\item Transitive verbs work differently: binary relations are from entities to events rather from entities to entities
				\item $discover\ phobos \Rightarrow \{(hall, \{event1045, event1046\})\}$
			\end{itemize}
			\pagebreak
			\item {\em Unified EV-FLMS (UEV-FLMS)}
			\begin{itemize}
				\item Based on EV-FLMS
				\item Solves two problems with how prepositions are treated in EV-FLMS
				\item Expands definition of prepositions such that multiple properties can be queried
				\item Novel method of handling the word ``by'', treating it as a preposition
				\item Unifies semantic concepts described in EV-FLMS
				%\item Improves efficiency
			\end{itemize}
		\end{itemize}
	\end{frame}
	
	\subsectionwithouttoc{Thesis Statement}
	\begin{frame}{Thesis Statement}
		By integrating a novel event-based denotational semantics with a parser constructed as an executable attribute grammar, it is possible to create a highly modular and extensible Natural Language Interface to the Semantic Web that supports the use of prepositional phrases in queries.
	\end{frame}
	
	\subsectionwithouttoc{Proof of Concept}
	\begin{frame}{Proof of Concept}
		We prove the Thesis by creating an online English query interface to a triplestore containing thousands of facts about the solar system\cite{Solarman:2016}.
		
		Some example queries that can be handled by this system include:
		
		\begin{itemize}
			\item ``when was something discovered at mt\_wilson''
			\item ``how was the thing that was discovered at flagstaff discovered''
			\item ``what was discovered in 1877 in us\_naval\_observatory''
			\item ``what planet is orbited by a moon that was discovered in 1684''
			\item ``which vacuumous moon that orbits jupiter was discovered by nicholson or hall with a telescope in 1938 in mt\_wilson or mt\_hopkins''
		\end{itemize}
		
		The interface can be accessed at: \url{http://speechweb2.cs.uwindsor.ca/solarman2/}
	\end{frame}
	
	\sectionwithouttoc{Demonstration}
	\begin{frame}{Demonstration}
		\begin{itemize}
			\item Two interfaces
			\begin{itemize}
				\item Natural Language Interface: queries are English sentences
				\item Direct Query Interface: queries are Haskell expressions, useful for learning how the semantics work
				\item Both interfaces compatible with screen readers for accessibility
			\end{itemize}
		\end{itemize}
	\end{frame}

	\begin{frame}{Visualization of NL query process}
		\centering
		``who discovered two moons in 1877 with a telescope''
		\[\Downarrow \qquad \Downarrow \qquad \Downarrow \]
		\[ \mathtt{who}\ \left(\mathtt{discovered'}\ \left(\mathtt{two}\ \mathtt{moons}\right)\ \left[ \mathtt{in'}\ \mathtt{1877},\ \mathtt{with}\ \left(\mathtt{a}\ \mathtt{telescope}\right) \right] \right)  \]
		\[\Downarrow \qquad \Downarrow \qquad \Downarrow \]
		\texttt{...(event1045, subject, hall), (event1045, year, 1877), (event1045, type, discover\_ev)...}
	\end{frame}

	%\sectionwithouttoc{Implementation Details}
	\begin{frame}{Implementation Details}
		\begin{itemize}
			\item Implemented in Haskell
			\item Sets are represented using lists, URIs are represented using \texttt{String} type
			\item Binary relations represented as association lists
			\item SafeHaskell is used in the Direct Query Interface to prevent arbitrary code execution
			\item Both interfaces are compatible with screen readers (for accessibility)
			%\item The monadic parser was joined with the monadic semantics in combination with a grammar
			\item Basic query fusion: memoization at query level
			\item Getts module provides a general extensible interface to triplestores using a typeclass
			\item Fast path for constructing FDBR with \texttt{condense} function
			\begin{itemize}
				\item Takes association list grouped according to the first element and constructs FDBR in O(n) time
			\end{itemize}
			\item Available on Hackage as the XSaiga package\cite{XSaiga:2016}
			%\item The actual query logic is contained in the Getts module, which provides interfaces to SPARQL triplestores and also list-based in-program triplestores.  A general typeclass is provided for future extensibility
		\end{itemize}
	\end{frame}

	\sectionwithouttoc{Unified EV-FLMS}
	\begin{frame}[allowframebreaks]{New semantics: Unified EV-FLMS}
		\begin{itemize}
			\item Formalizes notion of ``collecting'' binary relations using the term Function Defined By Relation (FDBR)
			\item Prepositions are denoted with a set of properties rather than a single property, necessary for ``in''
			\begin{itemize}
				\item i.e. $\left(\{props\}, predicate\right)$ instead of $\left(prop, predicate\right)$
			\end{itemize}
			\item Solved two problems with how prepositions were handled originally:
			\begin{itemize}
				\item Implicit ``and'' problem -- queries not honed down
				\item The problem of ``every'' -- predicates applied only to information in individual events
			\end{itemize}
			\item Unified semantic concepts
			\begin{itemize}
				\item Treating ``by'' as a preposition
				\item Uniform handling of passive and active verbs
				\item No bools: empty FDBR denotes False, non-empty FDBR denotes True -- nearly all semantic functions accept FDBRs and return FDBRs
				%\item Nearly all semantic functions accept FDBRs as arguments and return FDBRs (except question functions such as who, what, where, etc)
			\end{itemize}
		\end{itemize}
	\end{frame}

	\defverbatim[colored]\lstOFltEv{
	\begin{lstlisting}[basicstyle=\ttfamily,keywordstyle=\color{red}]
	filter_ev ev [] = True
	filter_ev ev (prep:list_of_preps)
		= ((snd (prep)) (getts_3 (ev,fst (prep),"?")))
		&& filter_ev ev list_of_preps
	\end{lstlisting}
	}

	\begin{frame}[allowframebreaks]{Problems with prepositions in EV-FLMS}
		Original Haskell \texttt{filter\_ev} function in EV-FLMS used for prepositions:
		{\lstOFltEv}
		\begin{itemize}
			\item The problem of ``every''
			\begin{itemize}
				\item Predicates are applied to the properties of individual events rather than to the entire set of events associated with a subject in an FDBR
				\item Fine for predicates such as ``a'', ``an'', and ``some''
				\item Determiners such as ``every'', ``two'' fail because they require information from all events
				\item Solution: collect all relevant properties from all events and invoke predicate on *those* instead
				\item Example: ``who discovered something with two telescopes''
				\item Two interpretations: discover one entity using two telescopes, or discover two entities using different telescopes
				\item Solution chosen allows for both interpretations
				\item In the future, users should be able to select which interpretation they want
				\item Clear-win example: ``who discovered with two telescopes'' 
			\end{itemize}
			\pagebreak
			\item Implicit ``and'' problem
			\begin{itemize}
				\item After solving the problem of ``every'', new problem: prepositions individually can't hone down events
				\item Result: implicit ``and'' placed between prepositional phrases
				\begin{itemize}
					\item ``in 1914 with two telescopes'' treated as ``in 1914 and with two telescopes'' -- subtle difference!
				\end{itemize}
				\item Re-thinking of problem required: predicates return booleans, but should return a new FDBR with only relevant FDBR-pairs
				\item Thinking along these lines resulted in the unification of several semantic concepts
			\end{itemize}

		\end{itemize}
	\end{frame}

	\begin{frame}[allowframebreaks]{Unified semantic concepts}
		\begin{itemize}
			\item All semantic functions except question functions (e.g. who, what, when, where...) accept FDBRs as input and return FDBRs
			\item Predicates return only FDBR-pairs in input that are {\em relevant} to that predicate
			\begin{itemize}
				\item An empty FDBR result denotes False, a nonempty FDBR result denotes True
				\item This is used in filtering according to chains of prepositions, solves both problems noted previously
			\end{itemize}
			\item Both passive and active transitive verbs handled uniformly
			\begin{itemize}
				\item Termphrase in active transitive verb, e.g. (a moon) in ``discover (a moon)'', is actually handled with new \texttt{filter\_ev}
				\item ``by'', as in ``discovered by'', is also handled with new \texttt{filter\_ev}
			\end{itemize}
		\end{itemize}
	\end{frame}
	
	%\sectionwithouttoc{Novel additions}
	\begin{frame}{Novel additions}
		The use of `by' as a preposition:
		\begin{itemize}
		\item In developing the system, it was discovered that the logic for handling the word `by' (e.g, ``phobos was discovered by hall'') was identical to a preposition in the semantics that searched for a subject in a given event
		\item The functionality was merged, and with no added complexity to the grammar, the word `by' is supported in queries
		\item Denotation for ``by'': $by\ termph\ = \left(\{``subject''\}, termph\right)$
		\end{itemize}
	\end{frame}
		
	\sectionwithouttoc{Timing}
	\begin{frame}{Timing}
		\begin{itemize}
			\item Two experiments
			\begin{itemize}
				\item Construct random binary relation with varying number of elements and compare run-time of \texttt{collect} implementations, including \texttt{condense}
				\item Perform two NL queries using an instrumented version of Solarman that collects profiling information
			\end{itemize}
			\item Experiment setup
			\begin{itemize}
				\item Intel Core i7 4770k Processor
				\item 16GB RAM 
				\item Samsung 850 EVO Solid State Drive
				\item 15 megabit cable connection
			\end{itemize}
		\end{itemize}
	\end{frame}

	\begin{frame}[allowframebreaks]{Experiment 1}
		\begin{itemize}
			\item Construct a randomized association list with 10,000 unique events and 1000 entities with varying number of pairs in the list
			\item Compare using both grouped and un-grouped association lists
			\item Compare performance of \texttt{collect} with previous implementation in \cite{agboola2015extensible}, and for grouped association lists, compare \texttt{condense} as well
			\pagebreak
			\item Results for ungrouped association lists:
			\begin{itemize}
				\item For 100,000 pairs:
				
				\begin{itemize}
					\item Previous implementation: 0.130 sec
					\item Our implementation: 0.129 sec
				\end{itemize}
				
				\item For 1,000,000 pairs:
				
				\begin{itemize}
					\item Previous implementation: 1.426 sec
					\item Our implementation: 1.388 sec
				\end{itemize}
				
				\item For 10,000,000 pairs:
				
				\begin{itemize}
					\item Previous implementation: 14.274 sec
					\item Our implementation: 14.283 sec
				\end{itemize}
			\end{itemize}
			\pagebreak
			\item Results for grouped association lists:
			\begin{itemize}
				\item For 100,000 pairs:
				\begin{itemize}
					\item Previous \texttt{collect}: 0.083 sec
					\item Our \texttt{collect}: 0.068 sec
					\item \texttt{condense}: 0.030 sec
				\end{itemize}
				
				\item For 1,000,000 pairs:
				\begin{itemize}
					\item Previous \texttt{collect}: 0.775 sec
					\item Our \texttt{collect}: 0.756 sec
					\item \texttt{condense}: 0.401 sec
				\end{itemize}
				
				\item For 10,000,000 pairs:
				\begin{itemize}
					\item Previous \texttt{collect}: 8.169 sec
					\item Our \texttt{collect}: 7.618 sec
					\item \texttt{condense}: 4.502 sec
				\end{itemize}
			\end{itemize}
		\end{itemize}
	\end{frame}

	\begin{frame}[allowframebreaks]{Experiment 2}
		\begin{itemize}
			\item Latency to remote SPARQL triplestore: approximately 60ms
			\item Two queries:
			\begin{enumerate}
				\item ``which vacuumous moon that orbits jupiter was discovered by nicholson or hall with a telescope in 1938 in mt\_wilson or mt\_hopkins''
				\item ``what was discovered in 1877 at us\_naval\_observatory''
			\end{enumerate}
			\item For each query, both IO and CPU time were measured using GHC's instrumentation
			\item Resultant ``.prof'' file produced by runs were examined to find where time was spent
			\pagebreak
			\item Results (profiling for IO):
			\begin{itemize}
				\item Query 1: approx 98.3\% of running time was spent waiting on IO
				\item Query 2: approx 98.3\% of running time was spent waiting on IO
			\end{itemize}
			\item Results (profiling CPU time, excluding IO time):
			\begin{itemize}
			\item Running time breakdown for query 1:
			
			\begin{itemize}
				\item SPARQL query generation and result processing including XML parsing: 51.4\%
				\item Network related processing: 39.3\%
				\item Natural Language parsing: 5.5\%
				\item Semantic functions: 3.8\%
			\end{itemize}
			
			\item Running time breakdown for query 2:
			
			\begin{itemize}
				\item SPARQL query generation and result processing including XML parsing: 42\%
				\item Network related processing: 44\%
				\item Natural Language parsing: 3\%
				\item Semantic functions: 11\%
			\end{itemize}
			\end{itemize}
		\end{itemize}
	\end{frame}

	\sectionwithouttoc{Future Work}
	\begin{frame}{Future Work}
		\begin{itemize}
			\item IO is clearly the limiting factor: many small queries produced, each query has overhead
			\item Scaling up:
			\begin{itemize}
				\item More advanced Query Fusion using monads
				\item Exploiting data-parallelism
				\item Foregoing set-theory and delving into conceptual spaces
			\end{itemize}
			\item Compatibility:
			\begin{itemize}
				\item Providing event based views of entity-based triplestores exploiting ontology data
			\end{itemize}
			\item Ambiguity:
			\begin{itemize}
				\item Provide a way for users to understand the various ways their queries may be interpreted and allow them to select which interpretation they mean
			\end{itemize}
		\end{itemize}
	\end{frame}

	\sectionwithouttoc{Conclusions}
	\begin{frame}[allowframebreaks]{Conclusions}
		\begin{itemize}
		\item Developed a novel event-based denotational semantics, UEV-FLMS, improving on EV-FLMS:
		\begin{itemize}
			\item Extends definition of prepositions to support multiple properties
			\item Solves implicit ``and'' problem and the problem of ``every'' in prepositions
			\item Novel method of handling the word ``by'', treating it directly as a preposition
			\item Unified semantic concepts
		\end{itemize}
		\item Implemented UEV-FLMS in Haskell--improved on efficiency of previous work
		\item Extended the work of Hafiz, Frost\cite{frosthafiz2008} to support non-referentially transparent functions as attributes
		\item Integrated UEV-FLMS with a parser constructed as an executable attribute grammar to produce an English Natural Language Interface with support for prepositional phrases in queries
		\item Developed an online interface to query facts about the Solar system as a demonstration
		\item Uploaded the work in the form of the XSaiga package on Hackage\cite{XSaiga:2016}
		\item Unifying semantic concepts may be a useful technique in solving other NL problems:
		\begin{itemize}
			\item Handling of Primary and Secondary sources in inference engines--notions of trust as an attribute
			\item Postpositions and circumpositions in other languages
		\end{itemize}
		\item There will be an increasing need for enabling technologies as IoT and Semantic Web gains mainstream usage
		\item Hopefully, the approach outlined here will prove fruitful for future researchers to build on in order to satisfy this need
		\end{itemize}
	\end{frame}

%	\begin{frame}{Future Publications}
%		\begin{itemize}
%			\item Nothing published yet
%			\item Want to submit to ESWC and ISWC conferences
%		\end{itemize}
%	\end{frame}
	
	\sectionwithouttoc{Thank You \& Feedback}
	\begin{frame}{Thank You \& Feedback}
		\centering
		Thank you for attending! \\ Comments, suggestions, or questions?
	\end{frame}
  
	\sectionwithouttoc{References}
	\begin{frame}[allowframebreaks]{References}
	  \def\newblock{}
	  \setbeamertemplate{bibliography item}[text]
	  \setlength\bibitemsep{1.5\itemsep}
	  \printbibliography
	\end{frame}
	
	%\appendix 
	%\section*{Appendix Overview}
	%\begin{frame}{Appendix Overview}
	%  \tableofcontents[part=3,sectionstyle=show/show,subsectionstyle=hide/hide/hide]
	%\end{frame}
	
	%\sectionwithouttoc{An Appendix Title}
	%\begin{frame}{\insertsection}
	%  This is the first section.
	%\end{frame}
	
	%\sectionwithouttoc{Another Appendix Title}
	%\begin{frame}{\insertsection}
	%  This is the second section.
	%\end{frame}
\end{document}
